% \subsection{\textcolor{red}{Поставленные задачи оценки и эталонные методы}}
\section{\textcolor{red}{Поставленные задачи оценки и эталонные методы}}

В рамках предложенного подхода оценки естественным образом
возникают две основные задачи:

\begin{enumerate}
    \item \textbf{Присвоение орбиталей.}
    Необходимо построить сегментацию пользователей на орбитали по их
    историческим признакам так, чтобы принадлежность к орбитали
    отражала устойчивые поведенческие паттерны, предсказывающие
    будущую выручку. Иными словами, орбиталь должна служить
    компактным суррогатом индивидуального поведения пользователя
    с точки зрения lLTV.

    \item \textbf{Оценка ценности орбиталей.}
    Требуется получить оценки инвариантных по времени множителей $\hat{V}_c$, 
    характеризующих относительную ценность каждой орбитали $c \in \mathcal{C}$ с точки
    зрения ожидаемой lLTV. Эти множители связывают
    поведенческие сегменты с их вкладом в будущую выручку.
\end{enumerate}

При наличии орбитальной принадлежности пользователей и набора
$\{\hat{V}_c\}_{c \in \mathcal{C}}$ можно вычислять метрики $r_V$ и
$r_{V\text{-corr}}$ по наблюдаемым распределениям пользователей по
орбиталям в начале и в конце эксперимента. Таким образом, относительный
эффект по lLTV оказывается выражен через изменения
распределения по орбиталям и их ценности, что позволяет оценивать
долгосрочное воздействие эксперимента, не дожидаясь полного завершения
прогнозного горизонта $\mathcal{F}(T_E, T_F)$.

Для оценки качества орбитального оценивателя рассмотрим два
естественных эталонных подхода.


\paragraph{(1) Прямое наблюдение ограниченной LTV.}

Наиболее прямой способ оценки долгосрочного эффекта состоит в том,
чтобы дождаться, пока прогнозное окно $\mathcal{F}(T_E, T_F)$ полностью
завершится, и затем непосредственно вычислить реализованную ограниченную
LTV каждого пользователя:
\[
  R_u(T_E) = \sum_{t' \in \mathcal{F}(T_E, T_F)} r_u(t').
\]
Далее средние значения $R_u(T_E)$ в контрольной и тестовой группах
используются для оценки относительного эффекта.


В идеализированной постановке такой оцениватель является несмещённым:
он опирается на фактические денежные потоки и не зависит от
модельных предположений. Однако на практике это свойство может
нарушаться, если за длительный период наблюдения поведение
пользователей в группах меняется по‑разному (например, из‑за
перекрёстных запусков других экспериментов, маркетинговых активностей
или сдвига состава аудитории). Кроме того, прямое наблюдение
ограниченной LTV страдает от высокой дисперсии, поскольку распределение
выручки по пользователям обычно имеет тяжёлые хвосты.


\paragraph{(2) Сильная предсказательная модель.}


Альтернативный эталонный подход состоит в построении
супервизируемой модели, предсказывающей $R_u(T_E)$ по признакам
пользователя, наблюдаемым к моменту $T_E$. Такая модель стремится
максимально использовать информацию о текущем состоянии пользователя
(частота и давность активности, структура трат, признаки вовлечённости
и т.п.), чтобы оценить его будущую ограниченную LTV.

Если модель хорошо специфицирована и обучена на репрезентативных
данных, она даёт оцениватель с существенно меньшей дисперсией по
сравнению с прямым наблюдением: индивидуальные шумовые колебания
денежных потоков сглаживаются за счёт агрегации информации в признаках.
Однако на практике такой подход чувствителен к ряду факторов:
ошибкам спецификации модели, сдвигу распределений признаков между
обучением и применением (covariate shift), переобучению и другим
источникам систематического смещения. В результате оценки эффекта
лечения могут быть искажены даже при визуально хорошем качестве
предсказаний.


\paragraph{Орбитальная оценка.}

Предлагаемый орбитальный оцениватель на основе метрик $r_V$ и
$r_{V\text{-corr}}$ занимает промежуточное положение между этими двумя
крайними подходами. С одной стороны, он опирается на структурное
модельное предположение
\[
  V_c(t) = V_c \cdot V_o(t),
\]
которое в реальных данных может выполняться лишь приближённо. Если
это предположение нарушается, оценки, основанные на орбиталях, могут
наследовать небольшое смещение.

С другой стороны, поведение пользователей сводится к относительно
небольшому числу орбиталей, и требуется оценить лишь низкоразмерный
вектор параметров $\{V_c\}_{c \in \mathcal{C}}$. Это радикально
снижает размерность задачи по сравнению с полностью индивидуальными
предсказаниями $R_u(T_E)$ и позволяет существенно уменьшить дисперсию
оценок эффекта. По сути, орбитали играют роль агрегированных суррогатов:
они достаточно грубо аппроксимируют сложное пространство
поведенческих траекторий, но при этом сохраняют различия между
типичными режимами поведения, важные с точки зрения будущей ценности.

Таким образом, орбитальный оцениватель реализует осознанный компромисс
между небольшим модельным смещением и существенным снижением дисперсии.
В контексте A/B‑экспериментов, где статистическая мощность часто
ограничивает возможность обнаружения слабых, но практически значимых
эффектов, такой компромисс оказывается особенно важен: снижение
дисперсии метрики непосредственно повышает чувствительность теста
при фиксированном объёме выборки.\textcolor{red}{[file:97]}


В \textcolor{red}{экспериментальной части работы} этот компромисс будет
квантифицирован на реальных данных: будут сопоставлены прямое
наблюдение ограниченной LTV, сильная предсказательная модель и
орбитальный оцениватель по смещению, дисперсии и способности
корректно выявлять эффект в A/A‑ и A/B‑экспериментах.
